\subsection{A3R: exclusion of every positive two-parameter Raynaud polarization}
\begin{lemma}[Connection bound; Proved]\label{a:connection-bound}
Let $C$ be a smooth projective connected curve over an algebraically closed
field of characteristic $p>0$. Let $L$ have positive degree, and let
$0\to\OO_C\to E\to L\to0$ be nonsplit with split Frobenius pullback.
For $0\leq n<p$ and any line bundle $M$ with
$\deg M<n\deg L$, we have $\Hom_C(\operatorname{Sym}^n E,M)=0$.
\end{lemma}
\begin{proof}
The splitting $F^*E=\OO_C\oplus L^p$ is unique since
$\Hom(L^p,\OO_C)=0$. Its canonical connection preserves $\OO_C$.
The second fundamental map $\beta:L^p\to\Omega_C^1$ is nonzero:
otherwise the splitting is horizontal and Cartier descent splits the
original extension. The precise descent source is Katz, \emph{Nilpotent
connections and the monodromy theorem}, Theorem~5.1, printed pp.~190--191 (opened in this run).
For a homomorphism $\varphi:\operatorname{Sym}^nE\to M$, its Frobenius
pullback is horizontal for the canonical connections. Its component on
$L^{pn}$ is zero by the degree inequality. On the symmetric power, the
connection carries the summand of exponent $j$ to the summand of exponent
$j-1$ with coefficient $j\beta$, in addition to terms of the same exponent.
If its components of exponents at least $j$ are zero, horizontality forces
$j\beta$ times its component of exponent $j-1$ to vanish. Since
$1\leq j\leq n<p$ and $\beta$ is nonzero, that component vanishes too.
Induction gives $F^*\varphi=0$, hence $\varphi=0$ by faithful flatness of
Frobenius on the smooth curve. For $n=0$ the assertion is simply that a
negative-degree line bundle has no sections.
\end{proof}

\begin{proposition}[A3R; Proved]\label{a:raynaud-all}
For the smooth Raynaud cyclic-cover surface of Zheng's construction,
with $\ell\mid p+1$, $L=N^\ell$ on its Tango base curve and
$\deg N>0$, let $\widetilde E$ be the positive ramification section.
For every $a,b>0$, the line bundle
$Z_{a,b}=\OO_X(a\widetilde E)\otimes\phi^*N^b$ cannot be both very
ample and have $H^1(X,Z_{a,b}^{-1})\ne0$.
\end{proposition}
\begin{proof}
Set $d=(p+1)/\ell$. For $a=\ell$, the finite-cover pushforward in
Zheng's Proposition~3.3 is
$\bigoplus_{i=0}^{\ell-1}\OO_{\PP(E)}(-id-1)\otimes\pi^*N^{ip-b}$.
Its direct image to $C$ is zero. Relative duality for the quotient
convention $\pi_*\OO_{\PP(E)}(1)=E$ therefore gives
\[
 H^1(X,Z_{\ell,b}^{-1})=
 \bigoplus_{i=1}^{\ell-1}
 H^0\!\left(C,(\operatorname{Sym}^{id-1}E)^\vee
                         \otimes N^{ip-b-\ell}\right).
\]
Here $n=id-1$ lies between zero and $p-d$, hence below $p$, and
\[
 \deg(N^{ip-b-\ell})-n\deg L=-(i+b)\deg N<0.
\]
The lemma makes every summand zero.
Next $\widetilde E^2=\deg N>0$, and $\phi|_{\widetilde E}$ is the
identity of $C$. Thus $Z_{a,b}^{-1}|_{\widetilde E}$ has negative
degree $-(a+b)\deg N$. The divisor exact sequence consequently
injects $H^1(X,Z_{a+1,b}^{-1})$ into $H^1(X,Z_{a,b}^{-1})$.
We obtain vanishing for every $a\geq\ell$ and $b>0$.
For $a<\ell$, restriction to the cuspidal fibre is
$\OO_F(a\infty)$ and Proposition~\ref{a:fibre} rules out very
ampleness. These two cases exhaust all positive $a,b$.
\end{proof}

This excludes the entire displayed two-parameter family. It is not a
nonexistence theorem for other line bundles or other surfaces.
The connection argument and all inequalities received an independent
targeted check in this run. The exact Cech computation covers1008 cases;
it serves as a check on this proof, not as its substitute.
